%% Items JBI collects separately from the manuscript at submission:
%% Highlights (3-5 bullets, at most 85 characters each) and the Statement of significance.
%% Build: latexmk -pdf submission_extras.tex
\documentclass[11pt,letterpaper]{article}
\usepackage[T1]{fontenc}
\usepackage[utf8]{inputenc}
\usepackage{lmodern}
\usepackage[margin=1in]{geometry}
\usepackage{booktabs,tabularx}
\usepackage{xcolor}
\definecolor{stneg}{HTML}{B4233C}
\newcommand{\todo}[1]{\textcolor{stneg}{[TODO: #1]}}
\pagestyle{empty}

\begin{document}

\section*{Highlights}
\begin{itemize}
\item One LLM request per paper yields gene--disease pairs with checkable evidence.
\item Four statuses separate supported, hedged, negative and absent relations.
\item Precision 0.894 on 419 BioRED abstracts; few contradictions of negative labels.
\item Re-querying missed pairs added no relations: the gap is an evidence standard.
\item A reason-code operating point raised held-out F1 from 0.542 to 0.627.
\end{itemize}

\section*{Statement of significance}
\todo{check the required form in the JBI guide for authors}

\medskip
\noindent
\begin{tabularx}{\linewidth}{@{}>{\bfseries}p{3.4cm}X@{}}
\toprule
Problem & Curators need gene--disease relations with evidence, status and source. Extraction tools return only positive pairs, mostly from abstracts.\\
What is already known & Supervised extractors reach high F1 on in-domain benchmarks such as BioRED. Zero-shot LLMs have been reported to lag in recall and consistency.\\
What this paper adds & A frozen decision card that gives an LLM high precision, checkable evidence and four-way statuses on three benchmarks and on full text. Evidence that much of the measured gap reflects differing annotation standards. An operating point that recovers recall on held-out data.\\
\bottomrule
\end{tabularx}

\end{document}
